\chapter{導来圏。}
\VerdierPhysicalPageAnchor{0174}%
\VerdierSemanticAnchor{ch3d}
\begingroup
\small
\setlength{\parskip}{0pt}
\setlength{\emergencystretch}{1.5em}
\setlength{\abovedisplayskip}{4pt plus 2pt minus 2pt}%
\setlength{\belowdisplayskip}{4pt plus 2pt minus 2pt}%

\VerdierSemanticAnchor{ch3d-1}
\section*{1. 導来圏の定義。}

\noindent\textit{この節では、特に断らない限り、$\mathcal A$ はアーベル圏を表す。}

\VerdierSemanticAnchor{ch3d-1-1}
\subsection*{1.1. 記法と用語。}

\VerdierSemanticAnchor{ch3d-1-1-1}
\noindent\textbf{1.1.1.}\quad
$\mathcal A$ を加法圏（必ずしもアーベル圏でなくてよい）とする。
$\operatorname{comp}^{b}(\mathcal A)$（それぞれ
$\operatorname{comp}^{+}(\mathcal A)$、
$\operatorname{comp}^{-}(\mathcal A)$）を、$|i|$ が十分大きいとき
$X^i=0$（それぞれ $i$ が十分小さいとき、$i$ が十分大きいとき）を満たす
複体 $X$ で定まる $\operatorname{comp}(\mathcal A)$ の充満部分圏とする。
$K^{b}(\mathcal A)$（それぞれ $K^{+}(\mathcal A)$、
$K^{-}(\mathcal A)$）を、$\operatorname{comp}^{b}(\mathcal A)$
（それぞれ $\operatorname{comp}^{+}(\mathcal A)$、
$\operatorname{comp}^{-}(\mathcal A)$）の対象で定まる $K(\mathcal A)$ の
充満部分圏とする。

\VerdierSemanticAnchor{ch3d-1-1-2}
\noindent\textbf{1.1.2.}\quad
次に $\mathcal A$ を（アーベル）圏とし、$*$ を活字記号
$b$、$+$、$-$、「空」のいずれかとする。
$\operatorname{comp}^{*,b}(\mathcal A)$（それぞれ
$\operatorname{comp}^{*,+}(\mathcal A)$、
$\operatorname{comp}^{*,-}(\mathcal A)$）を、
$\operatorname{comp}^{*}(\mathcal A)$ の対象であって、$|i|$ が大きいとき
$H^i(X)=0$（それぞれ $i$ が小さいとき、$i$ が大きいとき）を満たす
複体 $X$ で定まる $\operatorname{comp}^{*}(\mathcal A)$ の充満部分圏とする。
同様に、$K^{*,b}(\mathcal A)$（それぞれ $K^{*,+}(\mathcal A)$、
$K^{*,-}(\mathcal A)$）を、$K^{*}(\mathcal A)$ の対象であって、$|i|$
が大きいとき $H^i(X)=0$（それぞれ $i$ が小さいとき、$i$ が大きいとき）
を満たす複体 $X$ で定まる $K^{*}(\mathcal A)$ の充満部分圏とする。
このようにして $K(\mathcal A)$ の八つの部分圏を得る：

\VerdierSemanticAnchor{ch3d-1-1-2-diagram}
\begin{center}
  \resizebox{0.88\textwidth}{!}{%
    $\begin{tikzcd}[ampersand replacement=\&,column sep=3.5em,row sep=1.7em]
      \& \& K^{+}(\mathcal A) \arrow[dr] \& \& \\
      \& K^{+,b}(\mathcal A) \arrow[ur] \arrow[dr]
        \& \& K^{,+}(\mathcal A) \arrow[dr] \& \\
      K^{b}(\mathcal A) \arrow[ur] \arrow[dr]
        \& \& K^{,b}(\mathcal A) \arrow[ur] \arrow[dr]
        \& \& K(\mathcal A) \\
      \& K^{-,b}(\mathcal A) \arrow[ur] \arrow[dr]
        \& \& K^{,-}(\mathcal A) \arrow[ur] \& \\
      \& \& K^{-}(\mathcal A) \arrow[ur] \& \&
    \end{tikzcd}$%
  }\hfill\textup{(1.1.2.1)}
\end{center}
\endgroup
